\SourceLeaf{339}%
would then have followed at once. This makes no difference to \(\log\tau\), so we have \(\log t=2{,}6238578,\ t=420''\ 5889,\ T=5^0 38'51''4755\). Repeating the calculation with this value of \(T\) would leave \(B\) unchanged. For \(\log\tau\) one would find \(2{,}6238470\), giving \(\log t=2{,}6238573,\ t=420''\ 5884\) and hence \(T=5^0 38'51''4757\). Repeating the calculation once more with this value would give no change at all. Clearly, one could have stopped at the previous result, since with seven-figure logarithms the fourth decimal place of the second can still fluctuate by one or a few units. The final result is therefore
\[
\begin{gathered}
\text{Latitude of Inselsberg}=B^0-b=50^0\ 51'\ 8''\ 9444\\
\text{Azimuth of the side Inselsberg-Brocken}=180^0+T^0-t\\
=185^0\ 35'\ 21''\ 1815
\end{gathered}
\]
Finally, for the difference in longitude one finds
\[
\begin{aligned}
\log\lambda&=2{,}7313519\\
\log l&=2{,}7313472\\
l&=538''7002=0^0\ 8'\ 58''\ 7002
\end{aligned}
\]

There is no need to repeat the remarks made at the end of Article 20, which still apply in full here.
\par\medskip\centerline{\rule{.25\textwidth}{.3pt}}
